\documentclass[12pt]{article}
\usepackage[T1]{fontenc}
\usepackage[utf8]{inputenc}
\usepackage{lmodern}
\usepackage{textcomp}
\usepackage{enumitem}
\usepackage{geometry}
\geometry{margin=1in}
\begin{document}
\begin{center}
Answer Key - Economics - Graduate Level Assessment 24 \
\small (Answers are ordered exactly as in the question paper.)
\end{center}

\section*{Multiple Choice}
\begin{enumerate}[label=\arabic*.]
\item \textbf{Answer:} A
\\ \textit{Options:} A) increase government spending and taxes.; B) reduce government spending and increase taxes.; C) increase taxes without changing government spending.; D) maintain current levels of government spending and taxes.; E) reduce government spending and taxes.
\\ \textit{Note:} Increasing government spending and taxes during a recession stimulates aggregate demand, which can help boost economic activity and reduce unemployment, thereby aiding in recovery.
\item \textbf{Answer:} A
\\ \textit{Options:} A) The productivity of labor in country X is 75 percent higher than in country Z.; B) The productivity of labor in country X is 33 percent higher than in country Z.; C) The productivity of labor in country Z is 33 percent higher than in country X.; D) The productivity of labor in country Z is 50 percent higher than in country X.; E) The productivity of labor in country X is 50 percent higher than in country Z.
\\ \textit{Note:} The productivity of labor is calculated by dividing total output by the number of workers, so in country X, productivity is 10 units per worker (30,000 units/3,000 workers), while in country Z it is 5 units per worker (40,000 units/8,000 workers), making country X's productivity 100\% higher than country Z, not 75\%.
\item \textbf{Answer:} A
\\ \textit{Options:} A) The ability to buy goods and services; B) A financial instrument that represents ownership in a company; C) A system where goods and services are exchanged for commodities; D) A government-issued bond or security; E) A method of payment
\\ \textit{Note:} The correct answer is accurate because credit refers to a financial arrangement that allows individuals or businesses to obtain goods and services with the promise of future payment, thereby facilitating consumption and investment in the economy.
\item \textbf{Answer:} C
\\ \textit{Options:} A) They held the view that government intervention would reduce economic freedom.; B) They were convinced that laissez-faire would ensure equitable wealth distribution.; C) They were disturbed by the wide spread poverty and the fluctuations of the business cycle.; D) They thought laissez-faire would lead to more efficient and ethical business practices.; E) They favored the principles of socialism.
\\ \textit{Note:} The institutionalists of the early twentieth century rejected laissez-faire because they believed that unregulated markets contributed to social injustices, such as widespread poverty and economic instability, which necessitated a more interventionist approach to economic policy.
\item \textbf{Answer:} A
\\ \textit{Options:} A) no externalities and should not be subsidized; B) decreasing marginal utility and should be taxed; C) increasing marginal utility and should be subsidized; D) increasing marginal utility and should not be subsidized; E) externalities and should be taxed
\\ \textit{Note:} The correct answer reflects the concept that if education does not produce external benefits for society-such as increased civic participation or improved community outcomes-it undermines the justification for government subsidies, as the benefits are solely private and do not extend beyond the individual.
\end{enumerate}

\section*{True / False}
\begin{enumerate}[label=\arabic*.]
\setcounter{enumi}{5}
\item \textbf{Answer:} False
\\ \textit{Options:} A) True; B) False
\\ \textit{Note:} The statement is false because even if American agriculture is the most efficient, comparative advantage allows for beneficial trade, meaning the U.S. may choose to import certain foods that can be produced more cheaply or sustainably elsewhere.
\item \textbf{Answer:} False
\\ \textit{Options:} A) True; B) False
\\ \textit{Note:} The long-run Aggregate Supply (AS) curve is vertical, indicating that in the long run, output is determined by factors such as technology and resources rather than price levels or demand, meaning it does not slope upward with sustained increases in demand.
\item \textbf{Answer:} False
\\ \textit{Options:} A) True; B) False
\\ \textit{Note:} A depreciation of the dollar typically increases the cost of imports, making them less attractive to domestic consumers due to the higher prices of foreign goods.
\item \textbf{Answer:} False
\\ \textit{Options:} A) True; B) False
\\ \textit{Note:} The statement is false because an increase in disposable income typically leads to an increase in both consumption and saving, as individuals often allocate a portion of additional income to consumption rather than saving it all.
\item \textbf{Answer:} True
\\ \textit{Options:} A) True; B) False
\\ \textit{Note:} The statement is true because the equation represents a pure time series model where $y_{it}$ is influenced solely by temporal dynamics and a deterministic trend, without incorporating cross-sectional variations or additional predictors.
\end{enumerate}

\section*{Long Form Response}
\begin{enumerate}[label=\arabic*.]
\setcounter{enumi}{10}
\item \textbf{Answer:} Rejecting the null hypothesis in classical hypothesis testing indicates that there is sufficient statistical evidence to support the alternative hypothesis. However, it is incorrect to assume that the alternative hypothesis is automatically accepted; rather, it suggests that the null hypothesis is less plausible given the data. The alternative hypothesis may still be unproven or not practically significant, as rejecting the null does not confirm the alternative as true, but merely indicates that the evidence against the null is strong. Additionally, factors such as sample size, effect size, and the context of the study must be considered to evaluate the robustness of the findings. Hence, while rejecting the null hypothesis supports the alternative, it does not equate to its acceptance without further scrutiny and validation.
\\ \textit{Note:} correct because rejecting the null hypothesis signifies that the evidence favors the alternative hypothesis over the null, but it does not equate to accepting the alternative hypothesis as true or practically significant, as the rejection only indicates reduced plausibility of the null hypothesis based on the data.
\item \textbf{Answer:} In a monopolistically competitive industry, relatively free entry and exit are facilitated by several factors, including the role of long-term contractual obligations for suppliers. These obligations can create a barrier to entry if they restrict new entrants from securing necessary inputs or negotiating favorable terms. However, when contracts are flexible or designed with exit clauses, they enable suppliers to adapt to changing market conditions, thus reducing the risk associated with new market entrants. Additionally, low fixed costs and the availability of substitutes in monopolistic competition promote a dynamic environment where firms can enter and exit without significant financial repercussions. Consequently, the interplay between supplier contracts and market dynamics supports the fluidity of entry and exit in these industries.
\\ \textit{Note:} correct because it emphasizes that while long-term contractual obligations can create barriers to entry in a monopolistically competitive industry, flexible contracts with exit clauses allow suppliers to adapt and facilitate easier entry and exit for new firms, thereby promoting market dynamism.
\end{enumerate}

\vfill
\noindent\textit{End of Answer Key}
\end{document}